\documentclass[../main.tex]{subfiles}

\graphicspath{{../image/}{../figure/}}

\begin{document}

\chapter{前沿追踪法与标准 Riemann 半群}\label{chapter:ch14-前沿追踪与标准riemann半群}

前沿追踪方法提供了一种构造双曲守恒律系统 Cauchy 问题解的途径：用分段线性通量逼近非线性通量，从而把 Riemann 问题中的稀疏波替换为小强度激波（``稀疏前沿''）的组合，使波型由有限多个激波、稀疏前沿与人工伪激波组成。通过控制波相互作用并追踪其全局身份，可在整个上半平面上生成近似解，并在取极限后得到弱解。本章同时发展了标准 Riemann 半群的理论：在 $L^1(-\infty,\infty)$ 的闭子集上定义一族 Lipschitz 连续映射，其轨道即 Cauchy 问题的物理解，满足唯一性与连续依赖性。

\subfile{Chapter-14/section-01.tex}
\subfile{Chapter-14/section-02.tex}
\subfile{Chapter-14/section-03.tex}
\subfile{Chapter-14/section-04.tex}
\subfile{Chapter-14/section-05.tex}
\subfile{Chapter-14/section-06.tex}
\subfile{Chapter-14/section-07.tex}
\subfile{Chapter-14/section-08.tex}
\subfile{Chapter-14/section-09.tex}
\subfile{Chapter-14/section-10.tex}
\subfile{Chapter-14/section-11.tex}
\subfile{Chapter-14/section-12.tex}
\subfile{Chapter-14/section-13.tex}

\end{document}
